\section{The Turd in the Punch Bowl}

It will be useful to provide more concrete examples illustrating the difference between the consciousness formed by persuasion and that formed by rhetoric, as described by \textbf{Carlo Michelstaedter}.

\paragraph{The End of Rhetoric}


There is an intellectual lust just as surely as there is sexual lust. Men get excited about hearing a new idea, becoming part of a new movement, and engage in discussions, debates, and so on. Then the movement splinters into several factions. There is more rhetoric about how to name the factions. Then the debates devolve into contests between the factions and the aim of the original movement is pushed aside. Then it becomes personal: who is in and who is out. This is the path of infinite conversation and the way to resolve it.

Rhetoric $\Rightarrow$ endless conversation $\Rightarrow$ authority

Donoso Cortes\footnote{\url{https://gornahoor.net/?p=5016}} pointed out that rhetoric leads to infinite, and ultimately pointless, conversation. To end it, there has to be an ultimate authority. This is even more so for movements that oppose the amorphous, undifferentiated mass, since its opposite is organization and hierarchy. A movement cannot be in thought only, but must be alive and organized in the way that complements its thought.

Now ideally the man of authority is the one who knows. Of course, for spiritual movements, what he knows is the cosmic law. This law is supreme but adherents of movements make rhetoric itself supreme. Hence, they will accept members, even if they deny or oppose the cosmic law, provided the say and believe the right things.

Satan may throw great parties, but he always leaves his mark as a turd in the punch bowl. Those addicted to rhetoric will tell you to stop spoiling the fun. After all, the rest of it is all good and you can scoop around the turd if you are careful. I'm afraid no one is that careful.

\paragraph{Make your Paths Straight}


\begin{quotex}
Make your paths straight. 
\flright{\textsc{Matthew 3:3}}
\end{quotex}


The difference in understanding this, and similar, passages has changed considerably over the centuries. St Anthony of the Desert\footnote{\url{http://www.newadvent.org/cathen/01553d.htm}} explains:

\begin{quotex}
As far as the soul is concerned, being straight consists in its intellectual part's being according to nature, as it was created. ... Virtue exists when the soul maintains its intellectual part according to nature. It holds fast according to nature when it remains as it was made---and it was made beautiful and perfectly straight.
\end{quotex}


Hardly anyone today is concerned with the intellectual soul; almost no one can even recognize it in his own consciousness. Considering that it is the soul that is to be saved, this is quite strange, not to say, risky. The intellectual soul is straight when it is aligned with the cosmic law. Other words for straight are righteous or just. This indicates a persuaded man, in self-possession, indicating an inner, or esoteric, understanding.

Modern commentaries will interpret the passage as God making your path straight, or something similar. This is exoteric as it sees things coming from the outside, thus giving up one's own freedom and power.

\paragraph{Amor Fati}


\begin{quotex}
I have, in fact, brought this thing upon myself.
\flright{\textsc{Rudolf Steiner}}

There must be something in a human being that causes him to be born in a definite environment.  \flright{\textsc{Rudolf Steiner}}
\end{quotex}


According to Steiner, when something happens to a person that arouses a feeling of distress, he can take it in two different ways. The first way is to give into the negative feeling, and the second is to accept full responsibility. Now Steiner attributes the event to a past life, but that itself is a form of non-persuasion, pushing responsibility onto something we are not conscious of.

Nevertheless, the principle is sound. We should first look on our choices, or the string of choices, that led to the distress. We can often find the answer there. When we can't, that indicates a privation, a block, that we are not yet conscious of. Accept as much responsibility for your life as you can bear.

\paragraph{Consciousness of the Real I}


Someone mentioned today that his Buddhist teachers instructed him to be wary of the ``sense of I''. Apparently, he sees a conflict with the consciousness of the real I. This points out the danger of mixing traditions since vocabulary and frames of reference will be different.

The first principle is that everything must be verified. This is the opposite of rhetoric, since debates about various theories are not permitted. One of the things that we will verify that there is indeed a ``sense of I'', that is, I as an object. \emph{A fortiori}, there are many such I's.

Consciousness of the Real I is different as we can verify. It is not an object of sense, it is not even part of natural life. It is not subjective as are the I's of sense, but rather just and objective. Or ``straight'' if we prefer St Anthony the Great.

\paragraph{Pagan Unconscious}


\begin{quotex}
Our natural evolution in Western Europe was broken by the introduction of a psychology and of spirituality which had developed from a civilization higher than our own. We were interrupted at the very beginning when our beliefs were still barbarously polytheistic, and these beliefs were forced underground and have remained there for the last two thousand years. That explains the divisiveness found in the Western mind. Still in a primitive state, we were forced to adopt the comparatively sophisticated doctrines of Christian grace and love.
\flright{\textsc{Carl Jung}}
\end{quotex}


No longer able to contain such sophisticated doctrines, many Westerners today are returning to barbarously polytheistic beliefs. They are bitter against what they see as an alien imposition, both through physical force and effective propaganda. However, the return to the unconscious is the opposite of persuasion. It is made worse since the main adherents of the Christian religion no longer have any understanding of it.

They way forward is twofold. First, to understand the doctrines esoterically, which is really to give them life, beyond mere rhetoric. The second is to love our pagan past, using \textbf{Valentin Tomberg}, \textbf{Dante}, or \textbf{Boethius} as our models. This is an act of self-possession.

\paragraph{Gods within}


There is a pastime, especially popular among women, to take some test on the Internet that will reveal our true inner nature. Obviously, this involves giving up our power to some anonymous source. Our true inner nature can be revealed only through serious efforts at self-observation.

Jungian psychology to a large extent has devolved to a similar point. In this case, people want to know what ``inner god'' they are, and they take great pride in the answer. It is great party conversation in some circles. Unfortunately, that can only mean that one is possessed by unconscious forces represented by the god. This is the opposite of the self-possession that we want, or are.

As above, so below. Below, we integrate all these forces under the dominion of the True Self. Hence, above there is only one God.

\paragraph{Talking to Angels}


Another popular idea in some spiritual circles is the idea of communicating with angels. Once again, they are seen as ``outside'', and hence that involves surrendering our self-possession. This is safe, because there is the pleasure of the thought of angels, but no effort is involved on our part.

The persuaded man will understand the angels to be higher states of being, as Rene Guenon pointed out. So the angels must impart their own nature in order to be heard or understood.

\paragraph{Effects of Crystals}


\begin{quotex}
The divine sleeps in the rock, breathes in the plant, dreams in the animal, and wakes in man.
\flright{\textsc{Nordic adage}}
\end{quotex}


Since a crystal is a rock, it is part of the unconscious and hence represents a privation to us. Although Rene Guenon has taught us how to understand the conscious symbolism of highly developed traditions, Carl Jung has shown how the subconscious has its own symbols. Understanding such symbols can have a therapeutic effect. Hence, it is not surprising that crystals, through the symbolism associated with them, can also have healing effects.

As above, so below: any effects the crystals have come from the manifestation of more subtle powers. We can learn to possess them, too, and not project them onto the crystals.


\flrightit{Posted on 2014-06-11 by Cologero}

\begin{center}* * *\end{center}

\begin{footnotesize}
\begin{sffamily}
\texttt{Tom Blanchard on 2014-06-15 at 20:49 said:}

As usual, Cologero, you seem to have your finger right on the pulse of certain reactionary currents:
\url{http://www.radixjournal.com/blog/2014/6/15/wolves-among-the-ruins}


Vis a vis the ``return to the unconscious'' of paganism, I can't agree more. The bitterness felt by some is easy to sympathize with, because their problem is one that we also face. Nevertheless, the solution seems clear:

The problem was insoluble for pagan Icelanders because they had no conception of goods transcending ordinary natural purposes that could allow man's ultimate end --- the proper object of moral commitment --- to be distinguished from what particular men happen to want. Manly honor rooted in pure assertiveness was in the end self-contradictory, because in idealizing the assertive self it made what was asserted something ideal and therefore not merely the self. A man who asserted the self seemed to lack loftiness, while a more disinterested man seemed to lack energy; both qualities were necessary for true honor, but it seemed impossible to have them together. The problem was eventually resolved by adding Christian love to the native heroic conception to form a knightly ideal personified by Kari, Njal's son-in-law and avenger, who combined strength and courage with humility, forbearance, and in the end forgiveness.

(Taken from an article by Jim Kalb, written here: \url{http://antitechnocrat.net:8000/node/22} )


\end{sffamily}
\end{footnotesize}

